% ESMDA section.
% Tables are auto-generated: experiments_report/scripts/extract_esmda.py
%   -> sections/esmda_tables.tex (\EsmdaMatrixTable, \EsmdaMasterTable, \EsmdaCalibTable)
% Figures are auto-generated: experiments_report/scripts/make_esmda_figures.py
%   -> figures/esmda/

\begingroup\sloppy

\section{ESMDA}
\label{sec:esmda}

% AUTO-GENERATED by experiments_report/scripts/extract_esmda.py
% Do not edit by hand; re-run the script instead.
% 20 runs, 18 completed, 2 failed.

\newcommand{\EsmdaMatrixTable}{%
\begin{table}[htbp]
\centering
\caption{ESMDA: experiment matrix. Twenty configurations; the observation
interval is the method knob. \emph{loc} is localization on
(\code{loccorrelation}) or off (\code{locnone}). Run-directory names are given
here only; everywhere else runs are cited by ID. The two failed runs died in the
PALM \emph{truth} simulation before any assimilation started
(Section~\ref{sec:esmda-failed}).}
\label{tab:esmda-matrix}
\footnotesize
\begin{adjustbox}{max width=\textwidth}
\begin{tabular}{llllcc l}
\toprule
ID & run directory & truth & case & loc & obs [s] & status \\
\midrule
E1 & \code{pypalm\_to\_pyudales\_w3\_loccorrelation\_obs15\_inflow} & PALM & \code{inflow} & on & 15 & failed (truth run) \\
E2 & \code{pypalm\_to\_pyudales\_w3\_loccorrelation\_obs15\_inflow\_turb} & PALM & \code{inflow\_turb} & on & 15 & ok \\
E3 & \code{pypalm\_to\_pyudales\_w3\_loccorrelation\_obs15\_periodic} & PALM & \code{periodic} & on & 15 & ok \\
E4 & \code{pypalm\_to\_pyudales\_w3\_loccorrelation\_obs30\_inflow\_turb} & PALM & \code{inflow\_turb} & on & 30 & ok \\
E5 & \code{pypalm\_to\_pyudales\_w3\_loccorrelation\_obs30\_periodic} & PALM & \code{periodic} & on & 30 & ok \\
E6 & \code{pypalm\_to\_pyudales\_w3\_loccorrelation\_obs7p5\_inflow\_turb} & PALM & \code{inflow\_turb} & on & 7.5 & ok \\
E7 & \code{pypalm\_to\_pyudales\_w3\_loccorrelation\_obs7p5\_periodic} & PALM & \code{periodic} & on & 7.5 & ok \\
E8 & \code{pypalm\_to\_pyudales\_w3\_locnone\_obs15\_inflow} & PALM & \code{inflow} & off & 15 & failed (truth run) \\
E9 & \code{pypalm\_to\_pyudales\_w3\_locnone\_obs15\_inflow\_turb} & PALM & \code{inflow\_turb} & off & 15 & ok \\
E10 & \code{pypalm\_to\_pyudales\_w3\_locnone\_obs15\_periodic} & PALM & \code{periodic} & off & 15 & ok \\
\midrule
E11 & \code{pyudales\_to\_pyudales\_w3\_loccorrelation\_obs15\_inflow} & uDALES & \code{inflow} & on & 15 & ok \\
E12 & \code{pyudales\_to\_pyudales\_w3\_loccorrelation\_obs15\_inflow\_turb} & uDALES & \code{inflow\_turb} & on & 15 & ok \\
E13 & \code{pyudales\_to\_pyudales\_w3\_loccorrelation\_obs15\_periodic} & uDALES & \code{periodic} & on & 15 & ok \\
E14 & \code{pyudales\_to\_pyudales\_w3\_loccorrelation\_obs30\_inflow\_turb} & uDALES & \code{inflow\_turb} & on & 30 & ok \\
E15 & \code{pyudales\_to\_pyudales\_w3\_loccorrelation\_obs30\_periodic} & uDALES & \code{periodic} & on & 30 & ok \\
E16 & \code{pyudales\_to\_pyudales\_w3\_loccorrelation\_obs7p5\_inflow\_turb} & uDALES & \code{inflow\_turb} & on & 7.5 & ok \\
E17 & \code{pyudales\_to\_pyudales\_w3\_loccorrelation\_obs7p5\_periodic} & uDALES & \code{periodic} & on & 7.5 & ok \\
E18 & \code{pyudales\_to\_pyudales\_w3\_locnone\_obs15\_inflow} & uDALES & \code{inflow} & off & 15 & ok \\
E19 & \code{pyudales\_to\_pyudales\_w3\_locnone\_obs15\_inflow\_turb} & uDALES & \code{inflow\_turb} & off & 15 & ok \\
E20 & \code{pyudales\_to\_pyudales\_w3\_locnone\_obs15\_periodic} & uDALES & \code{periodic} & off & 15 & ok \\
\bottomrule
\end{tabular}
\end{adjustbox}
\end{table}%
}

\newcommand{\EsmdaMasterTable}{%
\begin{table}[htbp]
\centering
\caption{ESMDA: master results for all 20 runs, uDALES-truth block above,
PALM-truth block below. \emph{assim} / \emph{valid} are the vector RMSE
[m/s] of the ensemble mean at the 6 assimilated and 4 validation sensors,
averaged over time; \emph{field} $\Umag$ is the domain-wide RMSE of the velocity
magnitude. The two parameter columns give the knot-mean RMSE with the reduction
vs.\ prior in parentheses (priors are re-drawn per run, so only the reduction is
comparable across rows, and never across methods). ESMDA has no sequential
innovation, so $\chi^2$ is undefined; $z_{\mathrm{par}}$ is the pooled parameter
$z$-score standard deviation (1.03 = calibrated); $q$ is the field hit rate
(uDALES truth only). Bold = best in the column within the truth block; $q$ and
$\chi^2$ are not bolded.}
\label{tab:esmda-master}
\footnotesize
\begin{adjustbox}{max width=\textwidth}
\begin{tabular}{llcc rrr rr rrr r}
\toprule
 & & & obs & \multicolumn{2}{c}{vector RMSE [m/s]} & field & \multicolumn{2}{c}{parameter RMSE (red.\%)} & & & & wall \\
\cmidrule(lr){5-6}\cmidrule(lr){8-9}
ID & case & loc & [s] & assim & valid & $\Umag$ & $\alpha$ [deg] & $\Umag$ [m/s] & $\chi^2$ & $z_{\mathrm{par}}$ & $q$ & [h] \\
\midrule
E11 & \code{inflow} & on & 15 & 0.665 & 0.748 & 0.477 & 7.74 (61) & 0.529 (31) & -- & 12.1 & 0.72 & 2.19 \\
E18 & \code{inflow} & off & 15 & 0.694 & 0.754 & 0.521 & 7.59 (62) & 0.671 (18) & -- & 13.6 & 0.72 & 2.15 \\
E16 & \code{inflow\_turb} & on & 7.5 & 0.439 & 0.422 & 0.408 & 2.42 (86) & 0.379 (51) & -- & 6.1 & 0.74 & 2.24 \\
E12 & \code{inflow\_turb} & on & 15 & \textbf{0.401} & \textbf{0.365} & \textbf{0.372} & \textbf{2.31} (87) & \textbf{0.369} (53) & -- & 4.6 & 0.62 & 2.30 \\
E19 & \code{inflow\_turb} & off & 15 & 0.456 & 0.407 & 0.421 & 2.45 (87) & 0.439 (42) & -- & 7.8 & 0.62 & 2.29 \\
E14 & \code{inflow\_turb} & on & 30 & 0.418 & 0.366 & 0.388 & 2.35 (87) & 0.414 (48) & -- & 4.1 & 0.80 & 2.27 \\
E17 & \code{periodic} & on & 7.5 & 1.104 & 1.394 & 1.056 & 12.94 (18) & 0.754 (7) & -- & 3.3 & 0.84 & 1.73 \\
E13 & \code{periodic} & on & 15 & 1.091 & 1.408 & 1.040 & 13.62 (18) & 0.717 (9) & -- & 2.6 & 0.71 & 1.77 \\
E20 & \code{periodic} & off & 15 & 1.163 & 1.671 & 1.075 & 14.84 (22) & 0.764 (4) & -- & 89.8 & 0.58 & 1.80 \\
E15 & \code{periodic} & on & 30 & 1.089 & 1.417 & 1.033 & 13.63 (19) & 0.748 (7) & -- & 2.5 & 0.70 & 1.76 \\
\midrule
E1 & \code{inflow} & on & 15 & \multicolumn{9}{c}{\emph{failed (truth run)}} \\
E8 & \code{inflow} & off & 15 & \multicolumn{9}{c}{\emph{failed (truth run)}} \\
E6 & \code{inflow\_turb} & on & 7.5 & 1.432 & 1.327 & \textbf{1.024} & 7.83 (47) & \textbf{0.432} (41) & -- & 13.9 & -- & 2.24 \\
E2 & \code{inflow\_turb} & on & 15 & \textbf{1.422} & 1.377 & 1.082 & 7.12 (58) & 0.472 (40) & -- & 12.5 & -- & 2.36 \\
E9 & \code{inflow\_turb} & off & 15 & 1.481 & 1.389 & 1.050 & \textbf{6.62} (60) & 0.561 (30) & -- & 18.4 & -- & 2.35 \\
E4 & \code{inflow\_turb} & on & 30 & 1.429 & \textbf{1.323} & 1.068 & 7.44 (50) & 0.524 (35) & -- & 9.6 & -- & 2.30 \\
E7 & \code{periodic} & on & 7.5 & 1.495 & 1.806 & 4.033 & 14.86 (14) & 1.016 (-10) & -- & 3.2 & -- & 1.99 \\
E3 & \code{periodic} & on & 15 & 1.437 & 2.133 & 4.014 & 22.42 (-10) & 1.090 (-15) & -- & 3.1 & -- & 2.08 \\
E10 & \code{periodic} & off & 15 & 1.575 & 1.960 & 4.487 & 16.19 (2) & 1.602 (-36) & -- & 263.3 & -- & 2.17 \\
E5 & \code{periodic} & on & 30 & 1.438 & 1.941 & 4.190 & 15.86 (10) & 1.418 (-24) & -- & 3.2 & -- & 2.13 \\
\bottomrule
\end{tabular}
\end{adjustbox}
\end{table}%
}

\newcommand{\EsmdaCalibTable}{%
\begin{table}[htbp]
\centering
\caption{ESMDA: calibration diagnostics. $\chi^2$ is undefined for a smoother.
std$(z)$ is the standard deviation of the per-knot parameter $z$-scores
$(\theta^\ast-\bar\theta^a)/\sigma^a$; a calibrated 50-member ensemble gives
1.03, so every entry is over-confident. \emph{Contraction} is the knot-mean
posterior/prior spread ratio ($\ll 1$ = the data informed the parameter,
$\approx 1$ = unidentifiable). $q$ and its per-component values are field hit
rates (uDALES truth only). \emph{rank edge frac} is the fraction of the
validation-sensor rank histogram falling in the two outer bins; a calibrated
51-bin histogram gives 0.04.}
\label{tab:esmda-calib}
\footnotesize
\begin{adjustbox}{max width=\textwidth}
\begin{tabular}{llcc r rrr rr rrrr r}
\toprule
 & & & obs & & \multicolumn{3}{c}{std$(z)$} & \multicolumn{2}{c}{contraction} & \multicolumn{4}{c}{field hit rate} & rank edge \\
\cmidrule(lr){6-8}\cmidrule(lr){9-10}\cmidrule(lr){11-14}
ID & case & loc & [s] & $\chi^2$ & $\alpha$ & $\Umag$ & pool & $\alpha$ & $\Umag$ & $q$ & $u$ & $v$ & $w$ & frac (valid) \\
\midrule
E11 & \code{inflow} & on & 15 & -- & 12.9 & 11.5 & 12.1 & 0.30 & 0.37 & 0.72 & 0.93 & 0.73 & 0.50 & 0.50 \\
E18 & \code{inflow} & off & 15 & -- & 13.9 & 13.7 & 13.6 & 0.30 & 0.35 & 0.72 & 0.93 & 0.75 & 0.50 & 0.58 \\
E16 & \code{inflow\_turb} & on & 7.5 & -- & 5.2 & 7.0 & 6.1 & 0.23 & 0.36 & 0.74 & 0.94 & 0.65 & 0.61 & 0.75 \\
E12 & \code{inflow\_turb} & on & 15 & -- & 3.0 & 5.8 & 4.6 & 0.24 & 0.36 & 0.62 & 0.94 & 0.30 & 0.62 & 0.58 \\
E19 & \code{inflow\_turb} & off & 15 & -- & 4.5 & 10.2 & 7.8 & 0.22 & 0.29 & 0.62 & 0.94 & 0.30 & 0.62 & 0.33 \\
E14 & \code{inflow\_turb} & on & 30 & -- & 2.3 & 5.4 & 4.1 & 0.25 & 0.40 & 0.80 & 0.94 & 0.83 & 0.62 & 0.25 \\
E17 & \code{periodic} & on & 7.5 & -- & 3.0 & 3.6 & 3.3 & 0.59 & 0.71 & 0.84 & 0.93 & 0.88 & 0.70 & 0.17 \\
E13 & \code{periodic} & on & 15 & -- & 2.3 & 2.8 & 2.6 & 0.67 & 0.74 & 0.71 & 0.93 & 0.47 & 0.71 & 0.00 \\
E20 & \code{periodic} & off & 15 & -- & 91.9 & 88.8 & 89.8 & 0.19 & 0.19 & 0.58 & 0.88 & 0.21 & 0.64 & 0.17 \\
E15 & \code{periodic} & on & 30 & -- & 2.2 & 2.8 & 2.5 & 0.69 & 0.75 & 0.70 & 0.92 & 0.47 & 0.71 & 0.08 \\
\midrule
E1 & \code{inflow} & on & 15 & \multicolumn{11}{c}{\emph{failed (truth run)}} \\
E8 & \code{inflow} & off & 15 & \multicolumn{11}{c}{\emph{failed (truth run)}} \\
E6 & \code{inflow\_turb} & on & 7.5 & -- & 16.2 & 8.1 & 13.9 & 0.23 & 0.34 & -- & -- & -- & -- & 0.83 \\
E2 & \code{inflow\_turb} & on & 15 & -- & 13.3 & 7.0 & 12.5 & 0.24 & 0.37 & -- & -- & -- & -- & 0.83 \\
E9 & \code{inflow\_turb} & off & 15 & -- & 20.3 & 15.3 & 18.4 & 0.21 & 0.29 & -- & -- & -- & -- & 0.92 \\
E4 & \code{inflow\_turb} & on & 30 & -- & 10.1 & 6.3 & 9.6 & 0.25 & 0.39 & -- & -- & -- & -- & 0.67 \\
E7 & \code{periodic} & on & 7.5 & -- & 2.8 & 3.5 & 3.2 & 0.65 & 0.65 & -- & -- & -- & -- & 0.25 \\
E3 & \code{periodic} & on & 15 & -- & 3.3 & 2.8 & 3.1 & 0.64 & 0.80 & -- & -- & -- & -- & 0.25 \\
E10 & \code{periodic} & off & 15 & -- & 173.2 & 326.6 & 263.3 & 0.16 & 0.16 & -- & -- & -- & -- & 0.33 \\
E5 & \code{periodic} & on & 30 & -- & 3.0 & 3.3 & 3.2 & 0.65 & 0.82 & -- & -- & -- & -- & 0.33 \\
\bottomrule
\end{tabular}
\end{adjustbox}
\end{table}%
}


\subsection{Method and experiment matrix}
\label{sec:esmda-method}

ESMDA is run here as \code{TimeVaryingParameterESMDA} with
\code{joint\_state\_and\_parameter: false}, i.e.\ a \emph{parameter-only}
smoother. The state is never touched by the analysis: it improves only because
the forward model is re-run over the whole window with better boundary
parameters. Each of the three \SI{120}{\second} windows is solved with 3 ESMDA
iterations at inflation $\alpha_i = 3$ --- four forward ensemble passes per
window --- and is warm-started from the previous window's final state. The
control vector is the pair of time-varying inflow parameters $\alpha$ and
$\Umag$ on 15 harmonic knots (\SI{30}{\second} per knot, window boundaries
sharing a knot); \code{pin\_initial\_time\_point: true}, and the static vertical
inflow exponent is not estimated. Observations are time-mean aggregates of $u$
and $v$ over the observation interval, so the observation count per window is
$N_d = 192 / 96 / 48$ at $\Delta t_{\mathrm{obs}} = 7.5 / 15 / \SI{30}{\second}$.

Two knobs are swept. The \textbf{observation interval} is the method knob of
this section; localization on/off is the shared knob. Localization was varied
only at $\Delta t_{\mathrm{obs}} = \SI{15}{\second}$ and the interval only on
\code{inflow\_turb} and \code{periodic}, so the twenty configurations of
Table~\ref{tab:esmda-matrix} form a cross, not a full grid. The reference run of
each case is uDALES truth with localization on at
$\Delta t_{\mathrm{obs}} = \SI{15}{\second}$: E11 (\code{inflow}),
E12 (\code{inflow\_turb}), E13 (\code{periodic}).

One ESMDA-specific diagnostic is quoted throughout and appears in no table: the
normalised data mismatch $O_N$ (median over members) after each iteration, whose
nominal target is $1/2$. Every run is flagged \code{underfit\_final: true} with
\code{caveat: no\_representativeness\_error}, so only the \emph{trend} across
iterations is meaningful, and $O_N$ is not comparable between observation
intervals because $N_d$ differs.

\EsmdaMatrixTable

\subsection{Master results}
\label{sec:esmda-master}

\EsmdaMasterTable

Table~\ref{tab:esmda-master} collects every run. Read as a block it says four
things. First, the \emph{case} dominates every knob: with uDALES truth the
$\alpha$ RMSE is 2.31--\SI{2.45}{\degree} on \code{inflow\_turb},
7.59--\SI{7.74}{\degree} on \code{inflow} and 12.9--\SI{14.8}{\degree} on
\code{periodic}, while localization and the observation interval move it by
tenths of a degree. Second, that ordering \emph{inverts} the naive one: the
physically harder turbulent-inlet case is the one ESMDA does best on, because
there the boundary forcing --- and hence the estimated parameters --- is what
the sensors actually see. Third, the truth model matters almost as much as the
case: the same \code{inflow\_turb} configuration under PALM truth only reaches
6.62--\SI{7.83}{\degree}, and the field $\Umag$ RMSE degrades from
$\approx 0.37$ to $\approx 1.05$. Fourth, ESMDA \good{generalises at least as
well as it fits} wherever it works at all --- on E12 the validation vector RMSE
(0.365) is \emph{below} the assimilation vector RMSE (0.401), the expected
signature of correcting a global boundary condition rather than local state.

Two cautions on reading the table. The reduction-vs-prior percentages are
relative to each run's own re-drawn prior (\SI{14.6}{\degree} to
\SI{20.3}{\degree} of $\alpha$ RMSE across the set), so they compare rows within
this section but must never be carried across methods. And on \code{periodic}
the 18--22\,\% ``reduction'' is a spread reduction, not a bias reduction
(Section~\ref{sec:esmda-periodic}).

\subsection{Results by case}
\label{sec:esmda-cases}

This subsection uses uDALES truth only; PALM truth is
Section~\ref{sec:esmda-modelerror}.

\subsubsection{Laminar inflow (\code{inflow})}
\label{sec:esmda-inflow}

\begin{figure}[htbp]
\centering
\includegraphics[width=\textwidth]{esmda/valid_E11.png}
\caption{E11 (uDALES truth, \code{inflow}, localization on,
$\Delta t_{\mathrm{obs}} = \SI{15}{\second}$): validation sensors 0
($z = \SI{2}{\meter}$) and 3 ($z = \SI{20}{\meter}$, never observed at that
height) and the pooled validation error. The ensemble tracks truth almost
exactly for the first \SI{120}{\second}, then loses phase: the error rises from
$< 0.2$ to $1.7$ around $t \approx \SI{190}{\second}$ and the ensemble spread
does not cover the gap.}
\label{fig:esmda-valid-E11}
\end{figure}

\begin{figure}[htbp]
\centering
\includegraphics[width=0.92\textwidth]{esmda/perr_E11.png}
\caption{E11: per-knot parameter error, inflow angle above and velocity
magnitude below (RMSE of the ensemble mean and CRPS of the ensemble). The shaded
band is assimilation window~1. The angle error is small at the start and end of
the record and largest in the middle windows, mirroring the sensor phase error
of Figure~\ref{fig:esmda-valid-E11}.}
\label{fig:esmda-perr-E11}
\end{figure}

\begin{figure}[htbp]
\centering
\includegraphics[width=0.82\textwidth]{esmda/marg_E11.png}
\caption{E11: final-knot ($t = \SI{360}{\second}$) prior and posterior marginals
of $\alpha$ (truth \SI{19.0}{\degree}) and $\Umag$ (truth
\SI{8.09}{\meter\per\second}). The annotated $z$ is
$(\theta^\ast - \bar\theta^a)/\sigma^a$: the angle lands on truth ($z = 0.07$)
while the velocity magnitude over-corrects to $\approx
\SI{8.97}{\meter\per\second}$ ($z = -4.51$).}
\label{fig:esmda-marg-E11}
\end{figure}

\begin{itemize}[nosep,leftmargin=1.4em]
  \item \textbf{State.} Partial success. Assimilation vector RMSE 0.665,
        validation 0.748, field $\Umag$ RMSE 0.477. The failure mode is a
        \emph{phase} error, not an amplitude error: up to $t \approx
        \SI{120}{\second}$ both validation sensors sit on truth, after which the
        ensemble reproduces the right shape roughly \SI{20}{\second} late
        (Figure~\ref{fig:esmda-valid-E11}).
  \item \textbf{Parameters.} $\alpha$ is recovered well
        (\SI{7.74}{\degree} RMSE, 61\,\% reduction vs prior) and the final-knot
        posterior sits on truth (Figure~\ref{fig:esmda-marg-E11}), but $\Umag$
        over-corrects: \SI{0.529}{\meter\per\second} RMSE, 31\,\% reduction, and
        a final knot $+2.33$ prior-$\sigma$ above truth.
  \item \textbf{Where it stalls.} Window~0 fits essentially to the noise floor,
        but the pooled mismatch then flattens:
        $46.8 \to 21.0 \to 18.1 \to 17.5$, with the last two iterations barely
        moving. \bad{The problem is not the first window} --- it is that the
        warm-started state entering windows~1 and~2 carries an error that
        boundary parameters cannot repair. This is the honest motivation for a
        sequential state filter: here the parameters are already close to right
        and the state is not.
  \item \textbf{Calibration.} Badly over-confident: pooled parameter
        $z$-score std 12.1 against 1.03 expected, and half the validation-sensor
        rank mass in the two outer bins (edge fraction 0.50 against 0.04).
        Contraction is 0.30 ($\alpha$) / 0.37 ($\Umag$), so the data \emph{are}
        informative --- the ensemble simply shrinks too far.
\end{itemize}

\subsubsection{Turbulent inflow (\code{inflow\_turb})}
\label{sec:esmda-inflowturb}

\begin{figure}[htbp]
\centering
\includegraphics[width=\textwidth]{esmda/valid_E12.png}
\caption{E12 (uDALES truth, \code{inflow\_turb}, localization on,
$\Delta t_{\mathrm{obs}} = \SI{15}{\second}$): validation sensors 0 and 3 and
the pooled validation error. The elevated sensor~3 is tracked over the full
\SI{360}{\second} although every assimilated observation is at
$z = \SI{2}{\meter}$; the error stays near 0.25 with no drift.}
\label{fig:esmda-valid-E12}
\end{figure}

\begin{figure}[htbp]
\centering
\includegraphics[width=0.92\textwidth]{esmda/perr_E12.png}
\caption{E12: per-knot parameter error. The angle error stays between
\SI{0.8}{\degree} and \SI{4.1}{\degree} over all 15 knots and falls in the last
window, against a prior RMSE of \SI{17.8}{\degree}; compare the
\SI{1}{\degree}--\SI{15}{\degree} range of the PALM-truth runs in
Figure~\ref{fig:esmda-grid-palm}.}
\label{fig:esmda-perr-E12}
\end{figure}

\begin{figure}[htbp]
\centering
\includegraphics[width=0.82\textwidth]{esmda/marg_E12.png}
\caption{E12: final-knot marginals. The angle posterior is tight and centred on
truth ($z = -0.26$); the velocity magnitude is equally tight but centred
\SI{0.47}{\meter\per\second} high at \SI{8.56}{\meter\per\second}
($z = -4.60$) --- the systematic $\Umag$ bias described in the text.}
\label{fig:esmda-marg-E12}
\end{figure}

This is the case where ESMDA works, and the best run of the section.

\begin{itemize}[nosep,leftmargin=1.4em]
  \item \textbf{State.} Best in the uDALES block on every sensor metric:
        assimilation vector RMSE 0.401, validation \good{0.365}, field $\Umag$
        RMSE 0.372. Validation \emph{below} assimilation, and the never-observed
        $z = \SI{20}{\meter}$ sensor tracked over the full record
        (Figure~\ref{fig:esmda-valid-E12}).
  \item \textbf{Parameters.} $\alpha$ RMSE \SI{2.31}{\degree} from a
        \SI{17.8}{\degree} prior --- an \good{87\,\% reduction} (87\,\% on CRPS
        as well) --- and $\Umag$ RMSE 0.369 (53\,\%), both the best in the set.
        The mismatch is the only one in the section that marches monotonically
        towards its target, $33.7 \to 4.9 \to 2.6 \to 2.4$.
  \item \textbf{The one blemish is $\Umag$.} The final-knot posterior is
        $+1.43$ prior-$\sigma$ above truth (Figure~\ref{fig:esmda-marg-E12}),
        and this high bias is \textbf{systematic}: the final-knot normalised
        error is positive in every uDALES-truth run with an inlet
        ($+0.78$ to $+5.93$), while $\alpha$ is essentially unbiased
        ($-0.69$ to $+0.19$). The most plausible reading is that the ensemble
        compensates for unmodelled sub-grid and state error by raising the inlet
        speed.
  \item \textbf{Calibration.} Still over-confident --- pooled $z$-score std 4.6,
        validation rank edge fraction 0.58 --- but the best-calibrated
        \emph{informative} run of the section. Contraction 0.24 / 0.36.
  \item \textbf{Field skill is not uniform.} The pooled hit rate is 0.62, made
        of $u = 0.94$ but $v = 0.30$: the streamwise component is captured, the
        cross-stream one largely is not, even in the best run.
\end{itemize}

\subsubsection{Periodic boundaries (\code{periodic})}
\label{sec:esmda-periodic}

\begin{figure}[htbp]
\centering
\includegraphics[width=\textwidth]{esmda/valid_E13.png}
\caption{E13 (uDALES truth, \code{periodic}, localization on,
$\Delta t_{\mathrm{obs}} = \SI{15}{\second}$): validation sensors 0 and 3 and
the pooled validation error. The ensemble mean is a nearly flat, climatological
line against a chaotic truth; individual members are turbulent but mutually
decorrelated, so the mean averages them away while the spread stays wide.}
\label{fig:esmda-valid-E13}
\end{figure}

\begin{figure}[htbp]
\centering
\includegraphics[width=0.92\textwidth]{esmda/perr_E13.png}
\caption{E13: per-knot parameter error. Note the $y$ scale: the angle error is
an order of magnitude above Figure~\ref{fig:esmda-perr-E12} and shows no
downward trend within or across windows.}
\label{fig:esmda-perr-E13}
\end{figure}

\begin{figure}[htbp]
\centering
\includegraphics[width=0.82\textwidth]{esmda/marg_E13.png}
\caption{E13: final-knot marginals. The posterior has barely contracted, and
what movement there is goes the wrong way: the angle posterior sits near
$-\SI{8}{\degree}$ against a truth of \SI{19.0}{\degree} ($z = 4.45$).}
\label{fig:esmda-marg-E13}
\end{figure}

Here ESMDA does essentially nothing, and the reason is a \emph{parameterization}
failure rather than an algorithmic one: with periodic lateral boundaries there
is no inlet, so $\alpha$ and $\Umag$ no longer control what the sensors see. The
sensitivity of the observations to the estimated parameters is close to zero and
the ESMDA gain has nothing to work with.

\begin{itemize}[nosep,leftmargin=1.4em]
  \item \textbf{State.} Validation vector RMSE 1.408 and field $\Umag$ RMSE
        1.040, roughly $4\times$ and $3\times$ the E12 values. The ensemble mean
        is climatology (Figure~\ref{fig:esmda-valid-E13}); any state-updating
        method that beats climatology beats ESMDA here by default.
  \item \textbf{Parameters.} The mismatch is \bad{flat}:
        $37.3 \to 36.2 \to 36.7 \to 37.2$, numerically unchanged over three
        iterations (and equally flat at the other two intervals and without
        localization). Contraction stays at 0.67 / 0.74 against 0.24 / 0.36 on
        E12 --- the textbook signature of an unidentifiable parameter.
  \item \textbf{The parameters move the wrong way.} The nominal 18\,\% and
        9\,\% ``reductions'' in Table~\ref{tab:esmda-master} are spread
        reductions: the final-knot $\alpha$ posterior sits $-2.64$
        prior-$\sigma$ from truth (Figure~\ref{fig:esmda-marg-E13}), and $-3.10$
        / $-2.32$ at \SI{7.5}{\second} / \SI{30}{\second}.
  \item \textbf{Calibration is ``good'' for the wrong reason.} E13 has one of
        the lowest pooled $z$-score std values in the section (2.6, beaten only
        by the equally inert E15) and a validation rank edge fraction of 0.00.
        That is not skill: an ensemble that never
        updates keeps its prior spread, and a wide enough spread always covers
        the truth. Calibration statistics must be read together with the
        contraction ratio.
\end{itemize}

\subsection{Model error (PALM truth)}
\label{sec:esmda-modelerror}

\begin{figure}[htbp]
\centering
\includegraphics[width=0.63\textwidth]{esmda/valid_modelerror_E12_E2.png}
\caption{Matched versus model-error on the turbulent-inflow case: validation
sensors 0 and 3 for E12 (uDALES truth, above) and E2 (PALM truth, below), both
with localization on at $\Delta t_{\mathrm{obs}} = \SI{15}{\second}$. Against
PALM truth the uDALES ensemble produces a smooth, plausible, but largely
decorrelated signal: the validation error is roughly four times larger and no
longer shrinks where the ensemble is confident.}
\label{fig:esmda-modelerror}
\end{figure}

The eight completed PALM-truth runs are an internally consistent model-error
story, and they fail the way theory predicts.

\begin{itemize}[nosep,leftmargin=1.4em]
  \item \textbf{The mismatch stalls high and early.} On E2 the pooled median
        goes $50.1 \to 32.6 \to 30.2 \to 30.1$: iterations 2 and 3 buy almost
        nothing, and the posterior still misses the data by
        $\sqrt{2 O_N} \approx 7.8\,\sigma_{\mathrm{obs}}$ against
        $2.2\,\sigma_{\mathrm{obs}}$ on the matched E12. Without localization
        (E9) the mismatch \emph{increases} after the first iteration,
        $49.4 \to 32.4 \to 33.6 \to 35.0$.
  \item \textbf{The posterior becomes confidently wrong.} At the final knot the
        E2 $\alpha$ posterior is a tight violin near \SI{3.8}{\degree} against a
        truth of \SI{19.0}{\degree} ($z = 12.2$), with a contraction ratio of
        0.24 --- as confident as the matched run and much further from truth.
        $\alpha$ RMSE is \SI{7.12}{\degree} (58\,\% reduction against 87\,\%
        matched) and validation vector RMSE 1.377 against 0.365.
  \item \textbf{The sensor fit is insensitive to everything.} Across all eight
        PALM runs the assimilation vector RMSE is 1.42--1.58 regardless of case,
        interval or localization: once the forward model is structurally wrong,
        the achievable sensor fit is set by the model discrepancy, not by the
        assimilation configuration (Figure~\ref{fig:esmda-modelerror}).
  \item \textbf{Periodic plus model error is the worst corner of the set.} All
        four PALM/\code{periodic} runs \bad{degrade $\Umag$ relative to their
        prior} ($-10$ to $-36\,\%$), and E3 also degrades $\alpha$
        ($-10\,\%$, \SI{22.4}{\degree} RMSE) --- the single worst run anywhere
        in this section. Their field $\Umag$ RMSE (4.01--4.49) is four times the
        uDALES/\code{periodic} value, most of which is the truth-model
        difference rather than the assimilation.
  \item \textbf{No laminar PALM row exists} (Section~\ref{sec:esmda-failed}), so
        the model-error study covers only \code{inflow\_turb} and
        \code{periodic}.
\end{itemize}

\subsection{Sensitivities}
\label{sec:esmda-sensitivities}

\subsubsection{Localization}
\label{sec:esmda-localization}

\begin{figure}[htbp]
\centering
\includegraphics[width=0.6\textwidth]{esmda/marg_loc_E13_E20.png}
\caption{Final-knot marginals on \code{periodic} with localization on (E13,
above) and off (E20, below). Without localization the posterior collapses to a
sliver at \SI{7.46}{\meter\per\second}, $z = 43.3$ from truth: confident
nonsense produced by spurious correlations where the parameters carry no
signal.}
\label{fig:esmda-loc}
\end{figure}

Localization was varied only at $\Delta t_{\mathrm{obs}} = \SI{15}{\second}$,
giving five matched pairs (E11/E18, E12/E19, E13/E20 with uDALES truth; E2/E9,
E3/E10 with PALM truth). Its importance is \emph{opposite-signed} across cases:
a mild accuracy knob where the parameters are identifiable, a safety net where
they are not.

\begin{itemize}[nosep,leftmargin=1.4em]
  \item \textbf{\code{inflow\_turb} (E12 vs E19):} a small, consistent gain.
        $\alpha$ RMSE is unchanged (2.31 vs \SI{2.45}{\degree}), but $\Umag$
        RMSE improves $0.439 \to 0.369$, validation vector RMSE
        $0.407 \to 0.365$, and $\Umag$ calibration improves markedly
        (std$(z)$ $10.2 \to 5.8$).
  \item \textbf{\code{inflow} (E11 vs E18):} almost no effect on $\alpha$
        (7.74 vs \SI{7.59}{\degree}) but a large one on $\Umag$
        (0.529 vs 0.671; 31\,\% vs 18\,\% reduction). Without localization the
        final-knot $\Umag$ over-corrects to $+5.93$ prior-$\sigma$.
  \item \textbf{\code{periodic} (E13 vs E20):} localization is
        \good{load-bearing}. Without it the parameter $z$-score std explodes to
        \bad{91.9} ($\alpha$) and \bad{88.8} ($\Umag$) against 2.3 / 2.8, the
        contraction ratio drops to 0.19 from 0.67 / 0.74, and the final-knot
        $\Umag$ posterior is pinned at \SI{7.46}{\meter\per\second} with
        $z = 43.3$ (Figure~\ref{fig:esmda-loc}). Validation vector RMSE also
        degrades, $1.408 \to 1.671$, and the field hit rate falls
        $0.71 \to 0.58$.
  \item \textbf{PALM truth (E3 vs E10):} the same pattern, amplified --- on
        \code{periodic} the parameter $z$-score std goes to \bad{173}
        ($\alpha$) and \bad{327} ($\Umag$) without localization, against
        3.3 / 2.8 with it.
\end{itemize}

Localization is therefore best sold as \emph{insurance against spurious
correlations when the sensitivity is weak}, not as an accuracy knob.

\subsubsection{Observation interval}
\label{sec:esmda-obsinterval}

\begin{figure}[htbp]
\centering
\includegraphics[width=0.78\textwidth]{esmda/perr_obsinterval_E16_E12_E14.png}
\caption{Observation-interval sweep on \code{inflow\_turb} with uDALES truth and
localization on: per-knot inflow-angle error for E16 (\SI{7.5}{\second}),
E12 (\SI{15}{\second}) and E14 (\SI{30}{\second}). The knot-mean RMSE differs by
\SI{0.11}{\degree} across a fourfold change in data volume; the densest setting
is the worst of the three.}
\label{fig:esmda-obsinterval}
\end{figure}

The sweep exists for \code{inflow\_turb} and \code{periodic}, with localization
on, for both truth models. Because observations are time-mean aggregates, a
shorter interval means more --- but noisier and more correlated --- data.

\begin{itemize}[nosep,leftmargin=1.4em]
  \item \textbf{Accuracy is flat, and slightly worse when denser.} On
        \code{inflow\_turb} with uDALES truth, $\alpha$ RMSE is
        $2.42 / 2.31 / \SI{2.35}{\degree}$ and validation vector RMSE
        $0.422 / 0.365 / 0.366$ at $7.5 / 15 / \SI{30}{\second}$
        (Figure~\ref{fig:esmda-obsinterval}). \textbf{``More data'' is not a
        defensible claim here.}
  \item \textbf{Calibration improves monotonically with a longer interval.}
        $\alpha$ std$(z)$ falls $5.2 \to 3.0 \to 2.3$ and $\Umag$ std$(z)$
        $7.0 \to 5.8 \to 5.4$; the validation rank edge fraction falls
        $0.75 \to 0.58 \to 0.25$. \SI{30}{\second} is the best-calibrated
        setting and \SI{15}{\second} the lowest-error one.
  \item \textbf{\SI{7.5}{\second} is worst on both counts}, and its mismatch
        also stalls higher, $O_N = 4.8$ against 2.4 / 2.0 (values at different
        $N_d$, so only the stalling is meaningful). This is consistent with
        correlated aggregation error being effectively double-counted at short
        intervals.
  \item \textbf{On \code{periodic} the interval changes nothing:} $\alpha$ RMSE
        $12.9 / 13.6 / \SI{13.6}{\degree}$ and a flat mismatch at every
        interval. Under PALM truth it is equally flat
        (7.83 / 7.12 / \SI{7.44}{\degree} on \code{inflow\_turb}).
\end{itemize}

\subsection{Calibration}
\label{sec:esmda-calibration}

\begin{figure}[htbp]
\centering
\includegraphics[width=\textwidth]{esmda/rank_E11_E12_E13.png}
\caption{Rank of the truth within the posterior ensemble at the assimilated
(top) and validation (bottom) sensors, for the three reference runs E11, E12 and
E13. The dashed line is uniform and the grey band its $\pm 2\sigma$ range. E11
and E12 are strongly U-shaped --- the truth falls outside the ensemble far more
often than a calibrated ensemble allows, even in the best run of the section ---
whereas E13 looks flat only because its posterior never left the prior.}
\label{fig:esmda-rank}
\end{figure}

\EsmdaCalibTable

Table~\ref{tab:esmda-calib} gives the full diagnostics.

\paragraph{ESMDA is systematically over-confident.} A calibrated 50-member
ensemble has parameter $z$-score std 1.03 (\code{expected\_std} in every
summary). The smallest value anywhere in the section is 2.5 --- and it belongs
to E15, a periodic run whose posterior barely moved. Among the runs that
actually learn something, the best is E12 at 4.6, and the informative laminar
runs sit near 12--14. \bad{ESMDA should not be compared on RMSE alone.}

\paragraph{Rank histograms agree.} For E12 the assimilation panel puts 37 of 144
cases in the lowest rank bin and 49 in the highest where uniform would be 14
(Figure~\ref{fig:esmda-rank}). The validation rank edge fraction --- the share
of the 51-bin histogram in the two outer bins, 0.04 when calibrated --- is
0.25--0.75 across the informative uDALES runs and 0.67--0.92 on the
PALM-truth turbulent-inflow runs.

\paragraph{The ensemble never collapses at the member level.} All 18 completed
runs report 50 unique members in all three windows, with
\code{min\_over\_median\_pairwise} between 0.18 and 0.38. The over-confidence
lives in the posterior spread, not in member degeneracy, and no run triggered
the failure-resampling path.

\paragraph{Field skill is not uniform across components.}
\code{field\_metrics} is written only for uDALES-truth runs. The streamwise
component is captured well everywhere ($u$ hit rate 0.88--0.94 in all ten runs),
but the cross-stream component is highly variable: $v = 0.30$ on E12 and 0.21 on
E20, against 0.83 on E14 and 0.88 on E17. Pooled hit rates span 0.58--0.84. Do
not claim uniform field skill.

\subsection{Failed runs and cost}
\label{sec:esmda-failed}

Two of the twenty configured runs, E1 and E8 (PALM truth, \code{inflow},
$\Delta t_{\mathrm{obs}} = \SI{15}{\second}$, localization on and off), produced
only a \code{config.yaml} and a \code{hydra/} directory --- no
\code{run\_summary.yaml} and no figures. Both died in the \emph{truth} run,
before any assimilation started; the two logs are identical in substance:

\begin{quote}\small\ttfamily\raggedright
PALM produced only 52 of 180 expected 3D outputs (128 missing). PALM terminates
with exit 0 on its own numerical divergence, so this most likely means the run
diverged and stopped early. Refusing to pad 128 frames by repeating the last
one.
\end{quote}

\noindent
raised as \code{CalledProcessError} on \code{palm (urban\_run)} from
\code{pypalm/\allowbreak forward\_model.py:\allowbreak\ \_fit\_output\_window}. The distinguishing
setting is \code{inlet\_turbulence.enabled=false} together with
\code{initial\_seed=false}: PALM was asked to run an inflow/outflow case with no
inlet perturbation at all, and diverged after $\approx \SI{104}{\second}$ of the
\SI{120}{\second} window. The consequence is structural --- there is no
laminar-inflow row in the PALM-truth block, so a complete three-case
model-error comparison cannot be drawn.

\paragraph{Cost.} Wall time is 1.73--2.36\,h per run (last column of
Table~\ref{tab:esmda-master}), i.e.\ four forward ensemble passes per window and
three times the forward evaluations of a single-pass filter. The spread is
driven by the solver, not the algorithm: the uDALES/\code{periodic} runs are
cheapest (1.73--1.80\,h) simply because uDALES runs faster there, and the most
expensive run is E2 at 2.36\,h.

\subsection{Takeaways}
\label{sec:esmda-takeaways}

\begin{itemize}[leftmargin=1.6em]
  \item \textbf{Identifiability decides everything.} With a turbulent inlet and
        matched models ESMDA cuts $\alpha$ error by 87\,\%
        ($17.8 \to \SI{2.31}{\degree}$) using three analyses in
        \SI{360}{\second}; with a laminar inlet 61\,\%; with periodic boundaries
        the parameters stop being physical controls and the reduction collapses
        to an 18\,\% spread reduction.
  \item \textbf{The ordering inverts.} The physically harder turbulent-inlet
        case is the one ESMDA does best on, because there the boundary forcing
        is what the sensors actually see.
  \item \textbf{Generalisation beats fitting.} On E12 the validation vector
        RMSE (0.365) is below the assimilation one (0.401) and the never-observed
        \SI{20}{\meter} sensor is tracked. Correcting a global boundary condition
        carries no ``fits the observed sensors, fails elsewhere'' penalty.
  \item \textbf{Over-confident throughout.} The pooled parameter $z$-score std
        is 2.5--12.1 wherever localization is on, where 1.03 is calibrated; rank histograms are U-shaped
        everywhere, and the final-knot $\Umag$ posterior is above truth in every
        uDALES-truth run with an inlet, by $+0.78$ to $+5.93$ prior-$\sigma$.
  \item \textbf{Model error buys confidence.} Under PALM truth the
        mismatch stalls at $\approx 7.8\,\sigma_{\mathrm{obs}}$ instead of 2.2,
        the final-knot $\alpha$ posterior sits at \SI{3.8}{\degree} against truth
        \SI{19.0}{\degree} with $z = 12.2$, and the sensor fit saturates at
        1.42--1.58 regardless of any assimilation setting.
  \item \textbf{Knobs matter less.} Localization buys
        $\approx 10\,\%$ of validation vector RMSE where parameters are
        identifiable but prevents outright collapse where they are not
        ($z$-score std 92 or 327 without it). Denser observations do not help and
        slightly hurt: \SI{7.5}{\second} is the worst interval on both accuracy
        and calibration.
  \item \textbf{The state is missing.} On \code{inflow}
        the parameters are already close to right while the state drifts out of
        phase after \SI{120}{\second}, and on \code{periodic} the ensemble mean
        is climatology. Both are arguments for a method that updates the state,
        which is what the following sections test.
\end{itemize}

\subsection{Appendix: all runs}
\label{sec:esmda-appendix}

\begin{figure}[p]
\centering
\includegraphics[width=\textwidth,height=0.92\textheight,keepaspectratio]%
{esmda/perr_grid_udales.png}
\caption{Per-knot parameter error (inflow angle above, velocity magnitude below
in each tile) for all ten uDALES-truth runs, ordered by case and then by ID.
Each tile carries its own ID, case and knob setting. Note that the $y$ scales
differ by up to an order of magnitude between tiles; the shaded band is
assimilation window~1.}
\label{fig:esmda-grid-udales}
\end{figure}

\begin{figure}[p]
\centering
\includegraphics[width=\textwidth,height=0.92\textheight,keepaspectratio]%
{esmda/perr_grid_palm.png}
\caption{Per-knot parameter error for all eight PALM-truth runs, laid out as in
Figure~\ref{fig:esmda-grid-udales} (E1 and E8, the laminar-inflow pair, failed;
see Section~\ref{sec:esmda-failed}). Even on the turbulent-inflow runs the
inflow-angle error never falls below $\approx \SI{1}{\degree}$ and peaks near
\SI{15}{\degree}, against \SI{0.8}{\degree}--\SI{4.1}{\degree} in the matched
E12.}
\label{fig:esmda-grid-palm}
\end{figure}

\endgroup
